\section*{第 \ref{ch:b2:gratings} 章——多波干涉与光栅}

\begin{solution}{exo:b2:gratings:1}
$a = \qty{2}{\micro m}$, $\sin\theta_p = 0.275p$: \qty{16.0}{\degree}, \qty{33.4}{\degree}, \qty{55.6}{\degree};
无四级。一级可见范围为 \qty{11.5}{\degree} 至 \qty{20.5}{\degree}：
\qty{9}{\degree}。在 \qty{30}{\degree} 处，$\sin\theta = 0.5 + 0.275p$：$p = +1$（\qty{51}{\degree}）位于
一侧，$p = -1$ 至 $-5$（$13^\circ$、$-3^\circ$、$-19^\circ$、$-37^\circ$、$-61^\circ$）位于
另一侧。
\end{solution}

\begin{solution}{exo:b2:gratings:2}
$a = \qty{1.67}{\micro m}$。一级：$\theta = \qty{20.7}{\degree}$，$\Delta\theta = \Delta\lambda/a\cos\theta =
\qty{3.9e-4}{rad}$，\qty{0.19}{mm}；二级：$\theta = \qty{45}{\degree}$，$\qty{1.0e-3}{rad}$，
\qty{0.51}{mm}。像素尺寸为 \qty{50}{\micro m} 或更小。
\end{solution}

\begin{solution}{exo:b2:gratings:3}
$N = 12\,000$: $R = 12\,000$, $24\,000$, $36\,000$; $\Delta\lambda = 0.046$, $0.023$, \qty{0.015}{nm};
双线（\qty{0.6}{nm}）与 H/D 线对（\qty{0.18}{nm}，需 $R = 3600$）
都很容易分开。
\end{solution}

\begin{solution}{exo:b2:gratings:4}
CD：紫光 \qty{14.5}{\degree}，红光 \qty{25.9}{\degree}；红光至二级存在，紫光可达
四级：两道（部分重叠的）彩虹。DVD：紫光 \qty{32.7}{\degree}，红光
\qty{71}{\degree}；红光仅有一级：一道更宽的彩虹。
\end{solution}

\begin{solution}{exo:b2:gratings:5}
(a) 零点位于 $\varphi = 2\pi m/N$，$m = 1, \dots, N - 1$：共 $N - 1$ 个零点，
其间有 $N - 2$ 个极大值。(b) 当 $N\varphi/2 = 3\pi/2$ 时，$I \approx I_0/\sin^2(3\pi/2N) \approx I_0(2N/3\pi)^2$：
$4/9\pi^2 = 0.045$ 倍于 $N^2I_0$。(c) $N = 3$ 时仅有一个次极大，位于
$\varphi = \pi$ 处，高为 $I_0$，即主极大 $9I_0$ 的九分之一；$N = 6$ 时有四个
次极大，强度仅为百分之几。(d) $\sin^2x/x^2$ 的中央瓣
包含 $90\%$ 的能量。
\end{solution}

\begin{solution}{exo:b2:gratings:6}
(a) 需 $p > \lambda_{\min}/(\lambda_{\max} - \lambda_{\min}) = 1.33$，即自 $p = 2$ 起。(b) $\lambda/p$。(c)
\qty{25}{nm}；将棱镜与光栅交叉放置，可使各级一个叠在
另一个之上。(d) $\Delta\lambda = \lambda^2/2L$，不足一纳米——
比光栅的小数百倍。
\end{solution}

\begin{solution}{exo:b2:gratings:7}
(a) $\sin\beta = \lambda/2a = 0.3$，$\beta = \qty{17.5}{\degree}$。(b) 单个刻面的衍射
形成以其镜面方向为中心的宽瓣；落入其中的各级
获得光。(c) 二级（$2 \times \qty{500}{nm} = \qty{1}{\micro m}$，
同一角度处）。(d) $R = 2W\sin\beta/\lambda = 3.6 \times 10^5$。
\end{solution}

\begin{solution}{exo:b2:gratings:8}
(a) $\dd\lambda/\dd x = a\cos\theta/pf = \qty{3.1}{nm/mm}$：\qty{50}{\micro m} 对应 \qty{0.16}{nm}。(b)
$R = 30\,000$：\qty{0.018}{nm}。(c) 狭缝的影响大九倍；在 $s = \qty{6}{\micro m}$ 时二者相等。
(d) 狭缝处的衍射，以及探测器光能不足。
\end{solution}

\begin{solution}{exo:b2:gratings:9}
(a) $p = 2L/\lambda = 16\,700$；自由谱区 $c/2L = \qty{30}{GHz}$，$\lambda^2/2L = \qty{0.036}{nm}$。(b) $\mathcal F
= \pi\sqrt{0.95}/0.05 = 61$；\qty{0.5}{GHz}，\qty{0.6}{pm}。(c) $p\mathcal F = 10^6$，为
\qty{10}{cm} 光栅的四倍。(d) 其自由谱区
不足一纳米：各级会相互交叠，除非光已预先经过滤波
只保留其中某一级。
\end{solution}

\begin{solution}{exo:b2:gratings:10}
(a) $|\sin\theta - \sin\theta_0| \le 2$。(b) 两端光线的光程差为
$W(\sin\theta - \sin\theta_0)$，至多为 $2W$。(c) $2 \times 0.3/5 \times 10^{-7} = 1.2 \times 10^6$；$c/R
= \qty{250}{m/s}$。(d) 谱线中心可定位到其宽度的百分之一，
宽度；千条谱线平均后，误差按 $\sqrt{1000}$ 的比例下降，即达到米每秒的精度，
但需要近乎苛刻的稳定性。
\end{solution}

\begin{solution}{exo:b2:gratings:11}
(a) 由下层晶面反射的波多走 $2d\sin\theta$ 的路程。
(b) $\sin\theta = 0.273p$：\qty{15.8}{\degree}，\qty{33.1}{\degree}，\qty{55.0}{\degree}。(c) $\lambda \gg 2d$：
无角度满足该方程；玻璃没有规则晶面，因此
只形成漫射环——它是非晶体。(d) $\sin\theta \le 1$。
\end{solution}

\begin{solution}{exo:b2:gratings:12}
(a) 相邻单元的相位差由路径贡献 $2\pi a\sin\theta/\lambda$，
由馈电贡献 $-\psi$；主极大位于 $\varphi = 0$。(b) 第一零点位于
$\Delta\varphi = 2\pi/N$ 处：$\Delta\theta = \lambda/Na\cos\theta$。(c) 第二个主极大位于
$\varphi = \pm2\pi$ 处需 $|\sin\theta \pm \lambda/a| \le 1$：对每个指向角都不可能，
仅当 $a \le \lambda/2$ 时成立。(d) $Na = \qty{0.96}{m}$：\qty{1.8}{\degree}；$50$ 步，
\qty{50}{\micro s}。
\end{solution}

\begin{solution}{pb:b2:gratings:1}
\textbf{1.} $a = \qty{0.833}{\micro m}$; $N = 144\,000$.

\textbf{2.} $\sin\theta = \sin20^\circ - \lambda/a$（该级位于
法线的另一侧）：\qty{550}{nm} 对应 \qty{18.5}{\degree}，\qty{400}{nm} 对应 \qty{7.9}{\degree}，\qty{700}{nm} 对应
\qty{29.9}{\degree}：光谱范围为 \qty{22}{\degree}。

\textbf{3.} $1/a\cos\theta = \qty{1.27e-3}{rad/nm}$；$f_2\times$：\qty{1.27}{mm/nm}，即
\qty{0.79}{nm/mm}，每像素 \qty{0.012}{nm}。

\textbf{4.} $R = N = 144\,000$：\qty{0.0038}{nm}，为一个像素的三分之一。

\textbf{5.} \qty{400}{nm} 的二级光落到 \qty{800}{nm} 的一级光
本应到达之处，位于红光之外：可见光内无重叠；重叠将
始于 \qty{350}{nm}，而该波长可由玻璃滤光片去除。

\textbf{6.} \qty{22}{\degree} $\times$ \qty{1}{m} $= \qty{38}{cm}$：探测器覆盖一个区段；
转动光栅即可选出该区段。

\textbf{7.} $\lambda = 2a\sin17^\circ = \qty{490}{nm}$.

\textbf{8.} $0.8 \times 0.4 = 32\%$.

\textbf{9.} 角度分别为 \qty{26.5}{\degree}、\qty{14.0}{\degree}、\qty{10.3}{\degree}：相对 \qty{550}{nm} 位置的偏移量分别为 $+14$、$-8$ 和
\qty{-14}{cm}。

\textbf{10.} \qty{0.18}{nm} 相当于 $15$ 个像素：已分开。

\textbf{11.} 相距 $50$ 个像素，每条宽 $4$ 个像素。

\textbf{12.} \qty{100}{\micro m}，$6.7$ 个像素，\qty{0.08}{nm}——为
光栅极限的二十倍：受狭缝限制。

\textbf{13.} 像宽 \qty{97}{\micro m}：宽 \qty{50}{\micro m} 的狭缝
只通过一半的光。

\textbf{14.} 清晰度相同而光更少：光子噪声增大。

\textbf{15.} $\Delta\lambda = \lambda v/c = \qty{0.066}{nm}$，即 $5.5$ 个像素。

\textbf{16.} 地球的速度使每条谱线全年内最多移动 $\pm5.5$
个像素；由星历表可知，其精度远优于
\qty{100}{m/s} 所需的 $0.3\%$。

\textbf{17.} \qty{2.2e-5}{nm}，$0.002$ 个像素——远小于谱线宽度。
若将一条 \qty{7}{px} 谱线的质心定到其百分之一，即 \qty{0.07}{px}；千条谱线
合并给出 $\sqrt{1000}$：$0.002$ 个像素——恰与信号相当。真实
仪器还需更高的分辨本领与稳定性。

\textbf{18.} $\Delta\sin\theta = (\lambda/a) \times 10^{-5} = 6.6 \times 10^{-6}$, $\Delta\theta = \qty{7e-6}{rad}$,
\qty{7}{\micro m}：每开尔文约半个像素——对 \qty{10}{m/s} 须控温至毫开尔文量级。

\textbf{19.} 其谱线与待测谱线同一时刻记录在同一
像素上，可给出波长标尺并追踪每次漂移。

\textbf{20.} $a = \qty{3.33}{\micro m}$，$N = 36\,000$，角度相同（\qty{18.5}{\degree}），$R = 4N
= 144\,000$，色散相同，自由光谱范围 $\lambda/4 = \qty{137}{nm}$：分辨本领与
色散仅取决于 $W$ 与角度；自由光谱范围
缩小了。

\textbf{21.} 自由光谱范围 $\lambda^2/2L = \qty{0.15}{nm}$（$12$ 个像素）；$\mathcal F = 30$；峰
宽 \qty{0.005}{nm}：探测器上是一排尖锐的谱线，
构成波长标尺。

\textbf{22.} $2 \times 10^{-4}$，与 \qty{1.3e-3}{rad/nm} 相比：小六倍，并且
棱镜的分辨本领（$b\,\dd n/\dd\lambda \sim 10^4$）和非线性刻度
也随之失去；光栅还可用于紫外波段。

\textbf{23.} $R \le 2W/\lambda$：在 \qty{500}{nm} 处，宽一米者可达 $4 \times 10^6$——提高
分辨本领的途径只有增大宽度与角度。

\textbf{24.} $10^4 \times 0.012 \times 0.32 = 38$，即每像素每秒的光子数；$3800$
即 \qty{100}{s} 内的计数，信噪比约为 $\approx 60$。

\textbf{25.} 色散由 $p/a\cos\theta$ 与照相机焦距决定；
分辨本领由受照宽度与角度（$pN$）决定；
自由光谱范围由级次（$\lambda/p$）决定。
\end{solution}
